\section{Conclusion and Future Work}
\label{sec:conclusion}

In this work, we introduced \textbf{SPS-ZCA} (Single-Pass Sample-Wise Routing with Zero-Shot Centroid Alignment), a training-free and computationally efficient dynamic model-merging framework designed to address the critical real-world latency and robustness bottlenecks of edge deployments. 
By integrating \textit{Single-Pass Activation-Space Dynamic Blending} (SPS), our approach completely bypasses the sequential micro-batch partitioning required by state-of-the-art MBH systems, achieving a constant $O(1)$ backbone execution time and delivering a massive projected \textbf{3.90$\times$ analytical execution speedup} on highly heterogeneous batches (which is fully attainable physically under the proposed compiled hardware-software co-designed loop layout in Appendix A). 
Simultaneously, our \textit{Zero-Shot Centroid Alignment} (ZCA) leverages robust representation-space centroids rather than noisy classification heads to route inputs. Combined with Unit-Norm Calibration (UNC) and GMM-based coordinate density estimation, SPS-ZCA provides an exceptionally robust framework that recovers \textbf{100.0\%} of the theoretical Expert Ceiling and detects out-of-distribution queries with a \textbf{95.2\%} true positive rate.

From a systems-ML perspective, SPS-ZCA shifts the design paradigm away from complex, training-intensive routing modules to lightweight, geometrically grounded activation-blending operators that maximize physical hardware efficiency. 
There are several promising directions for future work. First, we plan to extend the SPS-ZCA framework to autoregressive Large Language Models (LLMs) and multi-modal foundation models, where dynamic expert switching of massive adapters (e.g., LoRA suites for reasoning, coding, and translation) is highly demanded but heavily bottlenecked by key-value (KV) cache management. Second, we aim to design hardware-specific compiler optimizations to execute the parallel low-rank matrix multiplications of SPS natively inside low-power Neural Processing Units (NPUs) and edge TPUs, unlocking sub-millisecond, multi-task dynamic serving on the next generation of consumer devices. Third, we plan to compile and benchmark our native fused C++ loop layout (detailed in Appendix A) directly on physical edge hardware (e.g., Raspberry Pi 4 or Nvidia Jetson) using ONNX Runtime CustomOps, and extend these physical serving benchmarks to larger autoregressive LLMs (such as LLaMA-3-8B) on actual mobile NPUs or edge TPUs to empirically demonstrate full physical wall-clock execution speedups in production. Fourth, we intend to conduct a comprehensive empirical evaluation of our fine-grained mitigations (Hierarchical Centroid Clustering and low-complexity Supervised Head Fine-Tuning) on fine-grained benchmarks (such as CUB-200 or Stanford Cars) to thoroughly validate dynamic model merging under extreme representational overlap. Fifth, we aim to standardize our custom C++ ONNX operators into mainstream edge compiler toolchains (such as Apache TVM, ExecuTorch, or MLC-LLM) to facilitate seamless, out-of-the-box integration for on-device modular serving.
